% TikZ 图：速度相关阻力与终端速度
% Day 06：牛顿第二定律扩展
\begin{tikzpicture}[>=Latex,scale=0.92,every node/.style={font=\scriptsize}]
  \begin{scope}
    \node[font=\small\bfseries] at (1.9,3.35) {受力变化};
    \draw[thick] (1.9,1.65) circle (0.28);
    \draw[blue,->,very thick] (1.72,1.35) -- (1.72,0.45) node[below] {$mg$};
    \draw[red,->,very thick] (2.08,1.95) -- (2.08,2.75);
    \node[red,right] at (2.08,2.55) {$f=kv$};
    \draw[->,thick] (2.75,1.95) -- (2.75,0.85) node[midway,right] {$v$};
  \end{scope}

  \begin{scope}[xshift=4.5cm]
    \node[font=\small\bfseries] at (2.1,3.15) {$v$-$t$ 图像};
    \draw[->,thick] (0,0) -- (4.25,0) node[right] {$t$};
    \draw[->,thick] (0,0) -- (0,2.85) node[above] {$v$};
    \draw[dashed] (0,2.15) -- (4.0,2.15);
    \node[left] at (0,2.15) {$v_t$};
    \draw[blue,very thick,domain=0:3.9,samples=80] plot(\x,{2.15*(1-exp(-0.9*\x))});
    \node[blue,align=center] at (2.55,1.20) {加速度逐渐减小};
    \node[align=center] at (2.55,2.55) {最终匀速\\$mg=kv_t$};
  \end{scope}
\end{tikzpicture}
